\chapter{Software realizzato}
\label{app:codice}

Questa appendice documenta l'organizzazione del software prodotto per la tesi e i
comandi con cui è stato usato. I listati completi non sono riportati integralmente:
si indicano struttura, scelte implementative rilevanti e modalità d'uso.

\section{Organizzazione del repository}

\begin{center}\small
\begin{tabularx}{\textwidth}{@{}lX@{}}
\toprule
\textbf{Percorso} & \textbf{Contenuto} \\
\midrule
\texttt{firmware/ld2410b\_logger/} & logger CSV a 5~Hz con engineering mode e
  lettura del PIR --- è il firmware usato per tutta la campagna \\
\texttt{firmware/test00\_ld2410\_info/} & lettura dei parametri del modulo via
  UART: firmware, MAC, soglie, gate, timeout \\
\texttt{firmware/test01\_ld2410\_base/} & sketch diagnostico con auto-scan del baud
  rate e stampa dei byte grezzi \\
\texttt{firmware/test03\_pir\_base/} & verifica del PIR e misura della durata degli
  impulsi \\
\texttt{firmware/test04\_set\_gate/} & configurazione del gate massimo per la prova
  di selettività spaziale \\
\texttt{firmware/ld2420\_monitor/} & lettura del LD2420 in modalità ASCII \\
\texttt{HLK-LD2410x/acquire.py} & acquisizione dalla seriale con metadati
  sperimentali (versione adattata) \\
\texttt{analisi/verifica\_engineering.py} & controllo di qualità dei file CSV \\
\texttt{analisi/analizza\_test.py} & metriche di confronto e aggregazione per
  scenario \\
\texttt{analisi/analizza\_respiro.py} & analisi in frequenza e ricerca del canale
  migliore \\
\bottomrule
\end{tabularx}
\end{center}

\section{Il logger}

\subsection{Scelte implementative rilevanti}

\begin{description}[leftmargin=1.4em,style=nextline,itemsep=3pt]
  \item[Periodo fisso e non ritardo fisso] Il ciclo principale calcola l'istante del
    prossimo campione anziché attendere un intervallo dopo l'elaborazione. È la
    ragione per cui la cadenza misurata è di 200~ms esatti su tutti gli intervalli,
    senza deriva: requisito necessario per l'analisi in frequenza.
  \item[Lettura del PIR nello stesso frame] Il valore del PIR è campionato
    immediatamente prima di comporre la riga CSV, quindi i due sensori condividono la
    base dei tempi per costruzione.
  \item[Engineering mode all'avvio] Il firmware entra in configurazione, attiva
    l'engineering mode, esce dalla configurazione e verifica che i frame ricevuti
    siano del tipo esteso, segnalando l'esito sulla seriale. Un fallimento silenzioso
    produrrebbe file senza colonne per-gate scoperti solo in fase di analisi.
  \item[Intestazione stampata una sola volta] L'intestazione CSV è stampata in
    \texttt{setup()}. Poiché il reset all'apertura della porta non è garantito, lo
    script di acquisizione la ricostruisce dal numero di colonne (si veda oltre).
\end{description}

\subsection{Formato di uscita}

\begin{lstlisting}[caption={Intestazione CSV completa in engineering mode, 29
colonne.},basicstyle=\ttfamily\scriptsize]
timestamp_ms,radar_presence,moving_target,stationary_target,
moving_distance_cm,stationary_distance_cm,moving_energy,stationary_energy,
pir_presence,
menergy_gate0,menergy_gate1,menergy_gate2,menergy_gate3,menergy_gate4,
menergy_gate5,menergy_gate6,menergy_gate7,menergy_gate8,
senergy_gate0,senergy_gate1,senergy_gate2,senergy_gate3,senergy_gate4,
senergy_gate5,senergy_gate6,senergy_gate7,senergy_gate8,
light_level,out_level
\end{lstlisting}

Le prime nove colonne coincidono con quelle del repository di
riferimento~\cite{repoCallisto}, per compatibilità.

\section{Acquisizione}

\subsection{Ambiente}

\begin{lstlisting}[language=bash,caption={Preparazione dell'ambiente Python.},
basicstyle=\ttfamily\scriptsize]
cd HLK-LD2410x
python -m venv .venv
.venv\Scripts\Activate.ps1
pip install pyserial numpy
mkdir data
\end{lstlisting}

\subsection{Uso}

\begin{lstlisting}[language=bash,caption={Esempio di acquisizione di un trial.},
basicstyle=\ttfamily\scriptsize]
python acquire.py --port COM3 --duration 300 ^
  --output data/fermo_seduto_T01.csv ^
  --scenario fermo_seduto --trial T01 ^
  --ground_truth_presence 1 --ground_truth_state still
\end{lstlisting}

Lo script aggiunge a ogni riga le colonne \texttt{pc\_time\_s},
\texttt{group\_id}, \texttt{trial\_id}, \texttt{scenario},
\texttt{ground\_truth\_presence}, \texttt{ground\_truth\_state}.

\subsection{La modifica necessaria}

La versione originale scrive righe solo dopo aver riconosciuto l'intestazione
stampata dall'ESP32 in \texttt{setup()}. Se il reset all'apertura della porta non
scatta, il file risulta \textbf{vuoto senza alcun messaggio}. La versione adattata
ricostruisce l'intestazione dal numero di colonne della prima riga di dati:

\begin{center}\small
\begin{tabular}{@{}rl@{}}
\toprule
\textbf{Colonne} & \textbf{Interpretazione} \\
\midrule
9 & modalità base \\
27 & engineering mode \\
29 & engineering mode con sensore di luce e stato OUT \\
\bottomrule
\end{tabular}
\end{center}

\section{Script di analisi}

\subsection{Controllo di qualità}

\begin{lstlisting}[language=bash,basicstyle=\ttfamily\scriptsize]
python analisi/verifica_engineering.py data/fermo_seduto_T01.csv
\end{lstlisting}

Verifica ed esce con diagnostica su: cadenza effettiva e jitter, presenza delle
colonne per-gate, percentuale di campioni saturi per canale, coerenza fra gate di
picco e distanza riportata, rumore di fondo per-gate nei tratti a stanza vuota,
numero e istante delle transizioni di presenza. Va eseguito su ogni file
\emph{prima} di usarlo per l'analisi.

\subsection{Metriche di confronto}

\begin{lstlisting}[language=bash,basicstyle=\ttfamily\scriptsize]
python analisi/analizza_test.py data/fermo_seduto_T0*.csv --salta-inizio 30
python analisi/analizza_test.py data/ingresso_T0*.csv --event-time 10
python analisi/analizza_test.py data/uscita_T0*.csv --release-time 10
\end{lstlisting}

Le opzioni principali: \texttt{-{}-salta-inizio} scarta i primi secondi (uscita
dell'operatore e stabilizzazione del PIR), \texttt{-{}-event-time} e
\texttt{-{}-release-time} attivano il calcolo delle latenze con la convenzione
dell'evento a istante noto. Con più file dello stesso scenario lo script aggrega e
riporta media e deviazione standard fra trial. Nel calcolo della latenza di rilascio
distingue i casi in cui il rilascio non avviene mai e quelli in cui avviene prima
dell'istante dichiarato, e conta le riaccensioni nella coda.

\subsection{Analisi del respiro}

\begin{lstlisting}[language=bash,basicstyle=\ttfamily\scriptsize]
python analisi/analizza_respiro.py data/fermo_seduto_T01.csv --scan
\end{lstlisting}

La modalità \texttt{-{}-scan} esamina tutti e venti i canali di energia, scarta
quelli saturi e quelli piatti secondo criteri dichiarati, ordina i rimanenti per
rapporto segnale-rumore del picco in banda 0,1--0,5~Hz e riporta la mediana delle
stime concordi con l'elenco dei canali che le sostengono. Esporta lo spettro in CSV
per la produzione dei grafici. Richiede NumPy; gli altri script usano la sola
libreria standard.

\section{Configurazione del gate massimo}

Lo sketch \texttt{firmware/test04\_set\_gate/} imposta il gate massimo via UART per
la prova di selettività spaziale. Tre accorgimenti sono degni di nota, perché
riguardano la scrittura di parametri persistenti su un modulo:

\begin{enumerate}[leftmargin=1.8em,itemsep=2pt]
  \item dopo ogni scrittura \textbf{rilegge sempre} i parametri per conferma, e
        segnala se il valore letto non corrisponde a quello scritto;
  \item \textbf{avvisa} quando la configurazione corrente è diversa da quella di
        fabbrica, per evitare di condurre acquisizioni con parametri modificati
        inconsapevolmente --- il rischio più insidioso in una campagna comparativa;
  \item il comando \texttt{d} riscrive i valori di fabbrica del nostro esemplare
        (le 9+9 soglie di \cref{tab:identita}, gate 8/8 e timeout di 5~s),
        così da poter tornare in ogni momento alla configurazione di riferimento.
\end{enumerate}
